\documentclass[a4paper,11pt]{article}
\usepackage[utf8]{inputenc}
\usepackage[T1]{fontenc}
\usepackage[french]{babel}
\usepackage{amsmath}
\usepackage[top=2cm, bottom=2cm, left=2.2cm, right=2.2cm]{geometry}
\usepackage{xcolor}
\usepackage{booktabs}
\usepackage{tabularx}
\usepackage{enumitem}
\usepackage{listings}
\usepackage{titlesec}
\usepackage{hyperref}
\hypersetup{colorlinks=true, linkcolor=primary, urlcolor=primary}

\definecolor{primary}{HTML}{1A237E}
\definecolor{success}{HTML}{2E7D32}
\definecolor{warning}{HTML}{E65100}
\definecolor{codebg}{HTML}{F5F5F5}

\lstset{backgroundcolor=\color{codebg}, basicstyle=\ttfamily\footnotesize,
  breaklines=true, frame=single, showstringspaces=false, columns=fullflexible,
  keywordstyle=\color{primary}}

\titleformat{\section}{\large\bfseries\color{primary}}{\thesection.}{0.6em}{}
\titlespacing*{\section}{0pt}{12pt}{4pt}
\titleformat{\subsection}{\bfseries}{\thesubsection.}{0.5em}{}
\setlength{\parindent}{0pt}
\newcommand{\code}[1]{\texttt{#1}}

\title{\textbf{Séance V2V --- Enrichissement et démonstration des messages CAM}}
\author{Douae Sebti}
\date{9 juillet 2026}

\begin{document}
\maketitle
\vspace{-8pt}\hrule\vspace{10pt}

% =====================================================================
\section{Objectif du jour}
Aujourd'hui, j'ai travaillé sur la communication \textbf{V2V} (\emph{vehicle-to-vehicle})
de ProtoVA, basée sur la pile \textbf{Vanetza} et les messages \textbf{CAM} de la norme
européenne ETSI C-ITS. Le CAM (\emph{Cooperative Awareness Message}) est le message
« je suis là » qu'un véhicule émet en continu (environ 1 fois par seconde) pour annoncer
sa position, son cap et sa vitesse aux véhicules voisins.

Mon objectif était double :
\begin{enumerate}[leftmargin=1.6em]
  \item \textbf{enrichir le CAM} pour qu'il porte le \textbf{cap} et la \textbf{vitesse}
        réels (et pas seulement la position) ;
  \item réaliser une \textbf{démonstration en direct} de l'échange des CAM entre les deux
        machines : le PC et la carte Jetson.
\end{enumerate}

% =====================================================================
\section{Ce que j'ai fait}

\subsection{Enrichissement du CAM (cap + vitesse)}
Jusqu'ici, le CAM émis contenait une position dynamique, mais le \textbf{cap} et la
\textbf{vitesse} étaient codés en dur à \code{0}. Je les ai rendus réels, en les faisant
remonter depuis l'odométrie du robot jusque dans le message. La chaîne complète est~:

\begin{lstlisting}[language=bash]
odometrie (nav_msgs/Odometry)
   |  cap = lacet (yaw) du quaternion   ;   vitesse = norme du twist
   v
odom_to_gps.py  ->  /v2v/my_gps = [lat, lon, cap, vitesse]
   v
v2v_node.py     ->  envoi UDP "lat,lon,cap,vitesse" a socktap
   v
socktap (Vanetza)  ->  remplit le CAM emis (Heading, Speed)
   v
WiFi  ->  autre vehicule  ->  CAM recu, decode
\end{lstlisting}

J'ai modifié cinq fichiers~:
\begin{itemize}[leftmargin=1.6em]
  \item \code{cam\_application.cpp} : le cap et la vitesse sont désormais lus depuis la
        position (\code{course}/\code{speed}) et convertis aux unités ETSI
        ($0{,}1^\circ$ pour le cap, $0{,}01$~m/s pour la vitesse). Si l'information n'est
        pas disponible, j'émets la valeur normalisée « \emph{unavailable} ».
  \item \code{udp\_position\_provider.cpp} : accepte maintenant un datagramme
        \code{"lat,lon,cap,vitesse"} (tout en restant compatible avec l'ancien
        \code{"lat,lon"}).
  \item \code{odom\_to\_gps.py} : calcule le cap (yaw) et la vitesse (twist) et les publie.
  \item \code{v2v\_node.py} : transmet ces deux champs à socktap.
  \item \code{v2v\_gps\_sim.py} : le véhicule simulé envoie aussi un cap et une vitesse
        cohérents.
\end{itemize}

J'ai recompilé \code{socktap} \textbf{sur le PC puis sur la Jetson} (même version de
Vanetza, commit identique), et j'ai vérifié qu'un envoi de
\code{48.7668,11.4320,90.0,5.00} produit bien un CAM avec \code{Heading 900} ($=90^\circ$)
et \code{Speed 500} ($=5$~m/s).

\subsection{La démonstration en direct (deux terminaux)}
J'ai lancé \code{socktap} \textbf{simultanément} sur les deux machines, chacune affichant
les CAM \emph{reçus} de l'autre~:

\begin{lstlisting}[language=bash]
# Terminal PC (station 1)
~/vanetza/build/bin/socktap -l udp -i wlp2s0 -a ca --print-rx-cam \
  --security none --station-id 1 -p static --latitude 48.7668 --longitude 11.4320

# Terminal Jetson (station 2)
~/vanetza/build/bin/socktap -l udp -i wlP1p1s0 -a ca --print-rx-cam \
  --security none --station-id 2 -p static --latitude 48.7669 --longitude 11.4321
\end{lstlisting}

Le lien \code{-l udp} utilise le multicast \code{239.118.122.97:8947} : les deux machines
étant sur le même WiFi, chacune reçoit les trames de l'autre, \textbf{sans sudo et sans
matériel radio 802.11p}.

% =====================================================================
\section{Résultats obtenus}
La démonstration a fonctionné dans les \textbf{deux sens} :
\begin{itemize}[leftmargin=1.6em]
  \item le \textbf{PC} reçoit les CAM du Jetson (\code{Station ID: 2},
        MAC \code{14:75:5b:15:59:32}) ;
  \item le \textbf{Jetson} reçoit les CAM du PC (\code{Station ID: 1},
        MAC \code{10:e1:8e:56:07:1a}).
\end{itemize}
Chaque CAM fait 99 octets et arrive environ une fois par seconde. Voici un CAM reçu,
décodé :

\begin{lstlisting}
CAM application received a packet with decodable content
Received CAM contains
  ITS PDU Header:  Protocol Version: 2  Message ID: 2  Station ID: 2
  CoopAwareness:
    Generation Delta Time: 58568
    Basic Container:
      Station Type: 5
      Reference Position:  Longitude: 114321000  Latitude: 487669000
      Altitude [Confidence]: 800001 [15]
    High Frequency Container [Basic Vehicle]:
      Heading [Confidence]: 3601 [10]
      Speed   [Confidence]: 16383 [1]
      Drive Direction: 0
      Vehicle Length: 1023   Vehicle Width: 62
received packet from 14:75:5b:15:59:32 (99 bytes)
\end{lstlisting}

\subsection{Signification de chaque champ}
\begin{tabularx}{\textwidth}{@{}l l X@{}}
\toprule
\textbf{Champ} & \textbf{Valeur} & \textbf{Signification} \\
\midrule
Protocol Version & 2 & version du format ETSI. \\
Message ID & 2 & \textbf{2 = CAM}. \\
Station ID & 1 / 2 & « plaque » de l'émetteur (PC=1, Jetson=2). \\
Generation Delta Time & 58568 & horodatage d'émission (ms). \\
Station Type & 5 & \textbf{voiture} (\emph{passenger car}). \\
Longitude & 114321000 & position en $10^{-7}$° $\rightarrow$ \textbf{11,4321°}. \\
Latitude & 487669000 & $\rightarrow$ \textbf{48,7669°}. \\
Semi Major/Minor Conf. & 500 & précision de la position : \textbf{5,00 m}. \\
Altitude [Conf.] & 800001 [15] & \textbf{indisponible} (pas de capteur d'altitude). \\
Heading [Conf.] & 3601 [10] & cap \textbf{indisponible} (voir remarque). \\
Speed [Conf.] & 16383 [1] & vitesse \textbf{indisponible} (voir remarque). \\
Drive Direction & 0 & marche avant. \\
Vehicle Length / Width & 1023 / 62 & dimensions \textbf{indisponibles}. \\
received packet from & MAC (99 o.) & adresse physique de l'émetteur + taille. \\
\bottomrule
\end{tabularx}

% =====================================================================
\section{Remarque importante : cap et vitesse « indisponibles »}
Dans cette démonstration, \code{Heading} vaut \code{3601} et \code{Speed} vaut
\code{16383}~: ce sont les valeurs \textbf{« unavailable »} normalisées par l'ETSI. Ce
n'est pas un bug, au contraire.

La raison est que j'ai lancé socktap en mode \code{-p static}, qui ne fournit
\textbf{que la position} (latitude/longitude) : aucun cap ni vitesse n'est donné. Mon code
enrichi réagit alors correctement en émettant la valeur « indisponible ». Cela
\textbf{prouve} que le repli fonctionne comme prévu.

Pour voir les \textbf{vraies valeurs} de cap et de vitesse dans le CAM, il faut utiliser le
mode dynamique \code{-p udp} et lui pousser \code{"lat,lon,cap,vitesse"} --- c'est
exactement ce que font \code{v2v\_node.py} et \code{odom\_to\_gps.py} à partir de
l'odométrie réelle du robot. J'ai vérifié ce mode séparément : le CAM affiche alors
\code{Heading 900} ($90^\circ$) et \code{Speed 500} ($5$~m/s).

% =====================================================================
\section{Bilan et prochaines étapes}
\textbf{Ce qui marche} : l'échange de CAM ETSI est opérationnel dans les deux sens entre le
PC et la Jetson, sur WiFi, avec position, identité et (en mode dynamique) cap et vitesse
réels issus de l'odométrie.

\textbf{À suivre} :
\begin{itemize}[leftmargin=1.6em]
  \item préparer une \textbf{démonstration CAM sur 5G} (attention : le multicast est en
        général bloqué par le réseau cellulaire, il faudra passer le lien socktap en
        \emph{unicast} ou par un relais) ;
  \item ajouter les messages \textbf{DENM} (alerte de freinage / d'obstacle) ;
  \item passer en \code{--security certs} pour des messages \textbf{signés} (vrai C-ITS).
\end{itemize}

\vspace{8pt}\hrule
{\footnotesize\itshape Documents liés : \code{v2v\_explique\_pas\_a\_pas.tex},
\code{doc\_v2v\_vanetza.tex}, \code{guide\_lancement\_v2v.tex}, \code{todo\_protova2.tex}.}

\end{document}
